\documentclass[12pt]{article}
\usepackage[utf8]{inputenc}
\usepackage[T1]{fontenc}
\usepackage{hyperref}
\title{Tai Loi Language Resources for K-12 Curriculum Development}
\author{Compiled by Claude Sonnet 5 (claude-sonnet-5) for Andr\'as Kornai}
\date{}
\begin{document}
\maketitle

\section{Introduction}
Tai Loi (ISO 639-3 tlq) is an Austroasiatic language of the Angkuic branch of Palaungic, spoken mainly in eastern Shan State, Myanmar (Joshua Project names the Kengtung, Mongping and Mongla townships), with smaller communities in Laos and China. Estimates are about 5,000 speakers, and the sources disagree on vitality (critically endangered, vulnerable or vigorous). ``Tai Loi'' is a Shan term meaning ``mountain Tai'', and Hall (2017) treats it as a cover term for several Angkuic, Waic and Western Palaungic languages rather than one language. Sources therefore disagree about what the label covers, and care is needed not to confuse it with Samtao or Mok.

\section{Primary Educational Resources}
The searches found no school materials in Tai Loi.

\section{Linguistic and Reference Materials}
\begin{itemize}
\item Wikipedia entry: \url{https://en.wikipedia.org/wiki/Tai_Loi_language}
\item Numerals data, Max Planck Institute for Evolutionary Anthropology: \url{https://lingweb.eva.mpg.de/channumerals/Tai-Loi.htm}
\item Numerals data for the related Samtao language: \url{https://lingweb.eva.mpg.de/channumerals/Samtao.htm}
\item Endangered Languages Project entry: \url{https://endangeredlanguages.com/lang/1245}
\item A paper from the University of Hawaii eVols repository that argues Tai Loi is not a single language (I have not read it): \url{https://evols.library.manoa.hawaii.edu/server/api/core/bitstreams/bf7256bd-c7cf-4a1d-bd93-b71253c4f8f7/content}
\end{itemize}

\section{Religious and Community Resources}
\begin{itemize}
\item Joshua Project people-group page for Tai Loi in Myanmar: \url{https://joshuaproject.net/index.php/people_groups/16146/BM}
\end{itemize}

\section{Digital and Audio Resources}
\begin{itemize}
\item Wikitongues language page: \url{https://wikitongues.org/languages/tlq/}
\end{itemize}

\section{Curriculum Development Context}
Little documentation exists, and the identity of the language is itself disputed. Any curriculum effort would first need to identify which communities and varieties are meant, then document them with the speakers.
\end{document}